\section{Introducción}

La densidad espectral de potencia (PSD) es una herramienta fundamental para analizar cómo se distribuye la potencia de una señal en frecuencia, lo cual permite evaluar ancho de banda, componente DC y ocupación espectral en sistemas de comunicaciones digitales \cite{ortega_reyes,proakis}. En señales aleatorias binarias, la PSD depende tanto de la estadística de la fuente como de la forma del pulso empleado en la transmisión.

En esta práctica se estudió la PSD de señales binarias aleatorias en GNU Radio Companion, considerando distintos valores de muestras por símbolo (\textit{Sps}), ruido blanco y fuentes reales obtenidas a partir de archivos de imagen y audio \cite{guia_lab3,gnuradio}. Además, se analizaron modificaciones en la codificación de línea y en la conformación de pulso para identificar su efecto sobre la componente DC y la distribución espectral de la señal.